\section{Heat-semigroup spectral barriers for nonzero-data unforcing}
\label{sec:xxx-spectral-barrier}

The forced velocity of Section~10 of \cite{openai2026} vanishes
initially. Edition~XXIX proved that a correction also starting from
zero cannot produce a nontrivial smooth unforced flow. We now develop
an independently checkable, frequency-resolved \emph{necessary}
condition for a correction starting from nonzero data. These estimates
are elementary Fourier--Duhamel consequences of the energy inequality,
not a new regularity theorem.

\begin{theorem}[Exact periodic Fourier energy barrier]
\label{thm:xxx-periodic}
Let $u$ be a smooth unforced three-dimensional incompressible
Navier--Stokes solution with viscosity $\nu>0$ on the torus
$\mathbb T^3=(\mathbb R/2\pi\mathbb Z)^3$. Normalize spatial volume
to one, put $E_0=\|u(0)\|_{L^2}^2$ and
$\widehat u_k=(2\pi)^{-3}\int_{\mathbb T^3}u(x)e^{-ik\cdot x}\,dx$.
For every nonzero $k\in\mathbb Z^3$ and $0\le t<T$,
\begin{equation}\label{eq:xxx-standard-mode}
\boxed{
|\widehat u_k(t)|\le e^{-\nu|k|^2t}|\widehat u_k(0)|+
\frac{1-e^{-\nu|k|^2t}}{\nu|k|}E_0.}
\end{equation}
If $\int_{\mathbb T^3}u(0)\,dx=0$, the stronger estimate holds:
\begin{equation}\label{eq:xxx-mean-zero-mode}
\boxed{|\widehat u_k(t)|\le
 e^{-\nu|k|^2t}|\widehat u_k(0)|+E_0G_{\nu,k}(t),}
\end{equation}
with
\[
G_{\nu,k}(t)=\begin{cases}
\displaystyle |k|\frac{e^{-2\nu t}-e^{-\nu|k|^2t}}
{\nu(|k|^2-2)},&|k|^2\ne2,\\[5pt]
\displaystyle |k|t e^{-2\nu t},&|k|^2=2.
\end{cases}
\]
Both formulas remain valid at all finite times before any alleged
singularity. They impose necessary conditions, not sufficient ones.
\end{theorem}
\begin{proof}
The projected Fourier equation reads
\[
\frac{d}{dt}\widehat u_k+\nu|k|^2\widehat u_k
=-i\mathbb P_k\sum_{\ell\in\mathbb Z^3}
 (k\cdot\widehat u_\ell)\widehat u_{k-\ell}.
\]
The right side has modulus at most
$|k|\sum_\ell|\widehat u_\ell||\widehat u_{k-\ell}|
\le |k|\sum_\ell|\widehat u_\ell|^2=|k|E(t)$,
by Cauchy--Schwarz and Parseval. The smooth energy identity gives
$E(t)\le E_0$. Integrating against the heat multiplier proves
\eqref{eq:xxx-standard-mode}. For mean-zero data, the mean stays
zero, Poincar\'e's inequality on the $2\pi$ torus gives
$E'(t)\le-2\nu E(t)$, hence $E(t)\le E_0e^{-2\nu t}$.
Integrating $|k|E_0e^{-\nu|k|^2(t-s)}e^{-2\nu s}$
proves the second formula, including the resonant
$|k|^2=2$ limit.
\end{proof}

\begin{theorem}[Whole-space dyadic shell and residual tradeoff]
\label{thm:xxx-whole-space}
Let $u=u_B+w$ be a smooth divergence-free finite-energy field on
$\mathbb R^3\times[0,T)$, with $u_B(0)=0$, $\nu>0$, and physical
residual
$R=\partial_tu+(u\cdot\nabla)u-\nu\Delta u+\nabla p$.
For $K>0$ let $\mathcal B_K$ be the $L^2$ Fourier projector onto
$\{\xi:K\le|\xi|\le2K\}$. Put
\[
x=\|w(0)\|_2,\quad \mathcal R_t=\int_0^t\|\mathbb P R(s)\|_2ds,
\quad b_K(t)=\|\mathcal B_Ku_B(t)\|_2,
\quad\delta_K(t)=\|\mathcal B_Kw(t)\|_2.
\]
Assume the standard smooth-solution energy identity and appropriate
spatial decay/integrability. Then
\begin{equation}\label{eq:xxx-physical-tradeoff}
\boxed{\quad
[b_K(t)-\delta_K(t)]_+\le
 e^{-\nu K^2t}x+
 \frac{\sqrt K}{\sqrt2\,\pi\nu}(x+\mathcal R_t)^2+
 \mathcal R_t.\quad}
\end{equation}
In particular, for an \emph{exactly unforced} completed flow
($\mathcal R_t=0$), writing
$a=e^{-\nu K^2t}$, $c=\sqrt K/(\sqrt2\pi\nu)$ and
$h=[b_K(t)-\delta_K(t)]_+$, a necessary explicit lower bound is
\begin{equation}\label{eq:xxx-nonzero-startup-floor}
\boxed{\quad
\|w(0)\|_2\ge
\frac{2h}{a+\sqrt{a^2+4ch}}\qquad(h>0).\quad}
\end{equation}
This lower bound applies directly to the authentic Section 10
background once $b_K$ has been sourced and independently certified.
It does not require pointwise evaluation of a whole-space Fourier
transform.
\end{theorem}
\begin{proof}
Use the unitary whole-space Fourier transform
$\widehat u(\xi)=(2\pi)^{-3/2}\int e^{-ix\cdot\xi}u(x)dx$.
For the projected quadratic nonlinearity,
\[
\left|\widehat{\mathbb P\operatorname{div}(u\otimes u)}(\xi)\right|
\le(2\pi)^{-3/2}|\xi|\|u\|_2^2,
\]
because $|\xi\cdot u|\,|u|\le|\xi|\,|u|^2$ under the
Fourier integral and $\|\mathbb P_\xi\|_{\mathrm{op}}=1$.
The projected residual equation and smooth energy inequality imply
$\sup_{s\le t}\|u(s)\|_2\le x+\mathcal R_t$.
Duhamel's formula, Minkowski, and Plancherel now yield
\[
\|\mathcal B_Ku(t)\|_2\le
 e^{-\nu K^2t}x+\mathcal R_t+
 \frac{(2\pi)^{-3/2}}\nu(x+\mathcal R_t)^2
 \left(\int_{K\le|\xi|\le2K}|\xi|^{-2}d\xi\right)^{1/2}.
\]
The shell integral is exactly $4\pi K$, giving
$\sqrt K/(\sqrt2\pi\nu)$. The triangle inequality
$\|\mathcal B_Ku_B\|_2\le\|\mathcal B_Ku\|_2+
\|\mathcal B_Kw\|_2$ proves
\eqref{eq:xxx-physical-tradeoff}. For $\mathcal R_t=0$,
solve $cx^2+ax-h\ge0$ for $x\ge0$ and rationalize the
positive root to obtain \eqref{eq:xxx-nonzero-startup-floor}.
\end{proof}

\begin{proposition}[Exact shear controls and limits of inference]
\label{prop:xxx-controls}
For each integer $n\ge1$, the periodic shear
$u=Ae^{-\nu n^2t}\sin(ny)e_x$ is an exact unforced
solution and obeys the purely linear part of the Fourier bound
with no quadratic transfer. In contrast
$u_B=t\sin(ny)e_x$ has projected physical force
$(1+\nu n^2t)\sin(ny)e_x$ and begins from zero; it is not an
unforced solution. Further, no pointwise whole-space bound of the
form $|\widehat u_0(\xi_0)|\le C\|u_0\|_2$ with a universal
constant $C$ is available for arbitrary Schwartz data; one must
use Fourier shells or a specified test function.
\end{proposition}
\begin{proof}
The shear fields are divergence-free and their convection term
vanishes identically. The heat equation gives the first field;
straight differentiation gives the stated force for the second.
For the Fourier pointwise claim, normalized smooth frequency bumps
supported in balls of radius $\varepsilon$ have unit $L^2$
norm but pointwise height of order $\varepsilon^{-3/2}$.
\end{proof}

\paragraph{Adversarial verification and limitations.}
An independently written exact SymPy replay checked the heat
convolution at $|k|^2=2$, the periodic shear equations, and the
unitary Fourier-shell constant. Forty thousand rational Cauchy
and energy-budget fixtures passed; a second independent audit
tested 1,200 divergence-free finite Fourier fields at four
wavevectors. This is a necessary condition; the constant in the
whole-space shell bound grows like $\sqrt K$, so it does not alone
exclude a shrinking-core singularity. No actual Fourier-shell
mass $b_K(t)$ or correction estimate $\delta_K(t)$ has yet been
certified for the complete authentic A.11/A.19-corrected source.
The end-to-end Navier--Stokes PDE and Lean formalization remain open.
